\section{Route Validation Rules}
\label{sec:interlocking:route-validation}

Route validation is the first and most critical step before granting authority for any train movement. It ensures that a requested path complies with all operational and safety constraints defined within the interlocking schema. To achieve this, the system evaluates the availability of every track segment along the route, confirming that no section is already occupied. It simultaneously checks for potential conflicts with active or pending paths to guarantee that no intersecting movements are permitted. The validation also extends to point machines, verifying that each one is correctly aligned and locked in its designated position. At the route’s endpoint, an additional safety requirement is enforced by reserving an overlap zone, effectively providing a braking buffer should the train fail to stop as expected.

All validation rules are maintained in station-specific JSON files that are parsed dynamically by the rule engine. This separation of rules from the codebase enhances flexibility, allowing adjustments or extensions without invasive software changes. In cases where any validation check fails, the route is immediately rejected, and the outcome is logged for full traceability. These logged results not only aid in diagnostic review but also form the foundation for subsequent signal aspect decisions, as detailed in Section~\ref{sec:interlocking:signal-decision}.